\chapter{Requirements under test}
\label{chap:product}

\section{Grouping the requirements}
Each of the 14 requirements was reviewed to identify the specified
behaviour, the inputs that could change the outcome, and the expected
result. Requirements that shared a decision or calculation were grouped
so that related scenarios could be designed together.

\begin{itemize}
    \item \textbf{Customer types and product categories (1--2).}
    These define the values used by the pricing rules.
    \item \textbf{Product discounts (3--4).}
    Both rules depend on quantities and on which entries count as the
    same product.
    \item \textbf{Coupons (5--8).}
    These rules jointly define which codes are accepted, which subtotal
    determines eligibility, and when the deduction is applied.
    \item \textbf{Monetary calculations (9, 10, and 14).}
    These specify the limits on deductions, how amounts are rounded,
    and how the final amount payable is calculated.
    \item \textbf{Delivery (11--13).}
    These rules use the delivery option, customer type, and discounted
    total to determine the fee.
\end{itemize}

The Java class structure guided the choice of where to test each rule.
For example, \texttt{DiscountService} implements both product discounts
and coupons. These were kept as separate requirement groups because their
outcomes depend on different decisions. Since one scenario could provide
evidence for several requirements, those links were recorded in the
requirements-to-tests table.

\section{Tracing scenarios to the source code}
After a scenario and its expected outcome were defined from the brief,
the responsible code was traced in three steps:

\begin{enumerate}
    \item The public method for the scenario was identified.
    \item Its calls were followed to track the values that affect the
    outcome.
    \item The class that makes the decision and the class that supplies
    its inputs were identified. This distinction guided the choice of a
    unit, integration, or mock-based test.
\end{enumerate}

For example, a Premium order that starts at \pounds78.00 and falls to
\pounds74.10 after a product discount should be charged \pounds4.99
for Standard delivery. The trace started at
\texttt{OrderService.createOrder} and followed the subtotal calculated by
\texttt{DiscountService}, the amount passed by \texttt{OrderService}, and
the comparison against \pounds75.00 in \texttt{DeliveryService}. The
threshold was tested directly in the delivery service, and the resulting
fee was checked in an integration test. Mockito was used to verify the
amount passed between the services.

